% !PRINTED 426
\weberpage{426}
Accordingly the formula now reads

\weq{10}{\Sigma k\, h(m) = 2\Sigma\partial,}

$n$ not a square,
\[
= 2\Sigma\partial + \sqrt{n} + \tfrac{1}{6},
\]
$n$ a square.

We decompose the sum on the left side of (10) into its individual parts by taking each solution of equation (8) separately. We then obtain a sum of expressions $h(m)$, in which, however, the same $m$ occurs as often as (8) has solutions. One of the solutions, however, gives
\[
m = \frac{4n-y^2}{x^2},
\]
and this occurs \gesperrt{once} if $y=0$, and twice (with $+y$ and $-y$) if $y$ is different from zero. But the number of primitive classes of the discriminant $-m$ is equal to the number of non-primitive classes of the discriminant $-x^2m$ of divisor $x$; and so, if we now understand by $h(m)$ (in distinction from the previous formula) the \gesperrt{total number} of classes of the discriminant $-m$ (primitive and imprimitive), then we obtain from (10):
\weq{11}{\begin{aligned}
&h(4n) + 2h(4n-1) + 2h(4n-4) + 2h(4n-9) + \cdots \\
&= 2\sum_{\partial>\sqrt{n}}\partial, \quad (n \text{ not a square}), \\
&= 2\sum_{\partial>\sqrt{n}}\partial + \sqrt{n} + \tfrac{1}{6}, \quad (n \text{ a square}).
\end{aligned}}

In this, the sum of the $h(4n-y^2)$ on the left is to be continued as long as $4n-y^2$ remains positive, and on the right $\partial$ runs through all divisors of $n$ that are greater than $\sqrt{n}$.

\gesperrt{It is, however, to be observed here that the forms $(x,0,x)$ are to be reckoned only with $\frac{1}{2}$ each, and $(x,x,x)$ with $\frac{1}{3}$,} because they arise from the representations of $4n$ by the discriminant $m=-4,-3$.

\begin{center}
\S\ 117. \textbf{\gesperrt{Arithmetic Nature of the Class Function $H_m(u)$.}}
\end{center}

We now return once more to equation (10), § 115, in order to derive from it the nature of the coefficients of the class equation $H_m(u)=0$. First, however, we consider the two special cases: $m=4$ and $m=3$.
